\documentclass[]{../../../../tex/classes/styledArticle}

\author{Shahar Perets}
\title{Podcast Script $\sim$ English Research Project $\sim$ COBE}
\begin{document}
	\maketitle
	\subsection*{Structure}
	\begin{itemize}
		\item Intro
		\item (I)
		\begin{itemize}
			\item Interview 3: with NYS Student
			\item Interview 1: with an Israeli student
			\item Interview 2: with a French Student
		\end{itemize}
		\item (II) Comparison
		\item (III) fake interview with Kish
		\item conclusion
	\end{itemize}
	
	The interviewer will be denoted in [I], and the interviewed person will be denoted in [II]. 
	
	\subsection*{Podcast}
	
	Hello, and welcome to the 42nd episode of ``History ou la mort''! Today, we have something special for you. Instead of doing another episode about historical events, we decided to review how students across the globe Learn History. We got tons of interviews for today, from all of the wide world, and even got some exclusive interviews with the Minister of Education in Israel, the glorified Mr. Yoav Kish. So let's dive in!
	
	Our first student for today is Samuel Barber, From New York, whom is going in two months to graduated from school. 
	\begin{itemize}
		\item[(I)] Samuel? 
		\item[(II)] I hear you well 
		\item[(I)] So, you're going to do the finals in history next week, right? 
		\item[(II)] what finals
		\item[(I)] In history..?
		\item[(II)] You meant social studies? 
		\item[(I)] Hmmm?
		\item[(II)] It's where we learn about things like culture, people, place, identity, technology... 
		\item[(I)] So you don't have a dedicated history lesson? Do you learn history at all? 
		\item[(II)] Yeah, we learn history through social studies. Everything from the agricultural revolution to the Second World War, and even the South African apartheid. 
		\item[(I)] Wow! That's extremely broad! How do you manage to learn the history 14$^{\text{th}}$ century, while learning to a great extent about events that happened in the 20th century, like when the Nazis fought the world war? 
		\item[(II)] Who are the Nazis? 
		\item[(I)] Seriously? 
		\item[(II)] I don't know, maybe we haven't gotten to it. 
		\item[(I)] You said you learned about The second WW
		\item[(II)] so..?
		\item[(I)] whatever
	\end{itemize}
	[end call]
	
	So, after this devastating beginning, we called, Francis Poulenc, a student from France. He graduated one year ago from a Technical School. 
	\begin{itemize}
		\item[(I)] Fancis? 
		\item[(II)] I hear you well
		\item[(I)] So, I'm doing a podcast episode about history learning in schools, and I thought you might be able to help me understand what you're learning over there in France. 
		\item[(II)] Oh, we learn a ton. Like everywhere else we learn about both world wars, the Holocaust and everything... We have some local history, before and after the second WW. For example, we dedicated many lessons to the French Revolution. 
		\item[(I)] Nice to know. It's important to learn modern local history too. wdym by ``The French Revolution''? It's broad subject. You learn everything from The Fall of the Bastille to the Reign of Terror and Napoleon? 
		\item[(II)] Terror? Anyway, yes. We learn how the Revolution began, and how Napoleon spread nationalist values to all of Europe. You know -- ``Liberté, égalité, fraternité''. 
		\item[(I)] There is a second part for this phrase, you know? Haven't you learn about Robespierre? 
		\item[(II)] who? 
		\item[(I)] The man... The Guillotine... Anyway, thanks for the call!
		\item[(II)] Bye bye! I'll research those. 
	\end{itemize}
	
	Our research extends beyond New York State and France, but in the end, we decided to interview an Israeli student, Paul Ben-Haim. 
	\begin{itemize}
		\item[(I)] Paul Ben Haim? 
		\item[(II)] I hear you well. I'm called Ben Chain. 
		\item[(I)] Ha- ha- Haim? 
		\item[(II)] Chaim
		\item[(I)] (coughing)aim? 
		\item[(II)] nvm
		\item[(I)] Alright, you're from Israel? 
		\item[(II)] yes, I am
		\item[(I)] Perhaps, you should guide us on what you are learning? 
		\item[(II)] On one hand, we cover general Nationalism, the events of world war 2 and world war 1... 
		\item[(I)] That isn't extremely broad. What about the Middle Ages, and the late 20th century?
		\item[(II)] The medieval period is covered in middle school, but no one remembers it because we're not tested on it. As for the 20th century, we learn about how we built a country in Israel. 
		\item[(I)] But what about cold war? The Armenian genocide? The rise and fall of Fascism and Communism? 
		\item[(II)] idk, some of those are mentioned briefly in the textbooks. What is Fascism and Communism? 
		\item[(I)] ok ok, so at least you can tell us what you learned about the Second World War? Any particular sections that really touched you? 
		\item[(II)] Fine. We learned last week about the youth and political movements during the holocaust. 
		\item[(I)] Like Aliyahat Hanoa'r, Hashomer Hatzair, The Bund --
		\item[(II)] The Bund? 
		\item[(I)] The socialist movement... The most popular among the jewish populations in the world war..?
		\item[(II)] We learn only about Zionist ones
		\item[(I)] Ah -- ooh -- I guess it's time to end the call, our time's up
		\item[(II)] ah, okay
	\end{itemize}
	
	After we gathered what was learned in history lessons throughout the globe (for our first research question), it's now time to understand why what's being learned, is learned. Almost everywhere, the local history is prioritized over general history, even if it means missing some important events. Apart from whatever is learned in New York, it seems like what is decided to be learned, is learned to a great extent -- for example, student spend a substantial part of the year learning the second world war, and not just 2 lessons. With that said, many times material that doesn't fit exactly the intentions of the state itself, is removed and not being learned. 
	
	To learn about how the history program is being decided in different countries, we got for interview the Mr. Yoav Kish. 
	
	\begin{itemize}
		\item[(I)] I have the privilege to talk to the minster of Education in Israel, Mr. Yoav Kish? 
		\item[(II)] (wait 3 sec)
		\item[(I)] I think you are on mute
		\item[(II)] Oh now you can hear me? 
		\item[(I)] I hear you well 
		\item[(II)] So... I'm doing a podcast about history lessons across the globe --
		\item[(I)] I know
		\item[(II)] -- ah, ikyk, my first question to you, is, how you decided what student learn and what they don't? 
		\item[(I)] So I'm Yoav, The great sun of my great father, Yakov, sincesinceibornbgan, --
		\item[(II)] I know
		\item[(I)] ikyk
		\item[(II)] Anyway, what I meant to say before I got interrupted, is that we teach our student the important values of country, nation, and tradition, and preserving the Jewish identity for the future. 
		\item[(I)] How does that relates to history
		\item[(II)] We teach them how the Zionist Nation spawn into existence, the heritage and values of the Jewish world... 
		\item[(I)] I think I got it
		\item[(II)] ok good I need to go to the Knesset destroy democracy I'm busy I need to go 
	\end{itemize}
	
	
	That's one way to end the call. Sometimes we can study more from what was taken out the syllabus, than what stayed in. For example, Israel decided to not teach about Socialist youth movements, since the Education Department would rather focus on the History of Zionism. Like Yoav Kish siad, focusing on nationalism and ignoring information that oppose that national movement, seems to be a way many countries choose. One would say this is crucial for the existence of the state, but other would critic them 
	
	for removing important information. 
	
	And now, we have got to our last research question -- how it impacts us? 
	
	First of all, it, as a way to understand the other. When talking and arguing with people who are not from Israel, it's important to bear in mind what the other side might know or not know, and vice versa. The majority of people do not know much more about history than what they're being taught in school; what is assumed as forgiven to us, could not be for the other, and being aware of the curriculum will help us add the required context. For example, it's useless to argue about Nazism, if the person whom you are talking with has never heard a formal definition of the Nazi ideology, and the only thing they know about the holocaust is that Hitler killed 6 million Jewish people (as is the case in New York). 
	
	Another way learning all of this can benefit us is by being a better student. Every book, test, and paper, especially in not-so-well-defined subjects like history, has the bias of its writers. And in the case of national education, all of that bias is governed by the state, which has its own intentions and biases. As students, by knowing what is expected of us, we can get a more informed and educated understanding of the history. 
	
\end{document}